%!TEX TS-program = xelatex
\documentclass[12pt,a4paper]{ctexart}

\usepackage{graphicx,booktabs,longtable,hyperref,multirow}
\usepackage[margin=2.3cm]{geometry}
\usepackage{titlesec,xcolor,fancyhdr,enumitem}

\hypersetup{colorlinks=true,linkcolor=blue,urlcolor=blue,
    pdftitle={第二阶段——问答数据集构建},pdfauthor={Qwen3-8B微调项目}}

\pagestyle{fancy}\fancyhf{}
\fancyhead[L]{Qwen3 微调项目}\fancyhead[R]{第二阶段报告}
\fancyfoot[C]{\thepage}

\title{\textbf{大模型微调项目第二阶段报告}\\[4pt]\large Raw Data $\to$ 指令数据集构建}
\author{计算机·密码与安全特化模型}
\date{\today}

\begin{document}\maketitle\thispagestyle{empty}

\begin{abstract}
\noindent 本报告记录了第二阶段的核心工作——将第一阶段产出的19,464个原始文本块（raw data）通过强模型辅助生成转化为高质量指令微调数据集。采用DeepSeek- chat模型作为生成引擎，设计了覆盖五种认知层次的Prompt模板，经过17批次API批量生成、质量过滤和分层抽样，最终产出3,678对结构化问答样本，按80:10:10划分训练集、验证集和测试集，可直接用于Qwen3-4B的QLoRA微调训练。
\end{abstract}

\section{阶段概述}

\subsection{目标}

在第一阶段原始文本采集和清洗的基础上，构建面向计算机科学、密码学与网络安全领域的高质量指令微调数据集。主要目标包括：

\begin{enumerate}[nosep]
    \item 从raw text chunks生成多样化问答对
    \item 覆盖概念解释、章节总结、方法对比、论文摘要理解和推理分析五种任务
    \item 建立自动质量过滤和人工审核机制
    \item 按80:10:10比例划分训练/验证/测试集
    \item 产出可直接用于HuggingFace QLoRA训练的JSONL数据集
\end{enumerate}

\subsection{输入与输出}

\begin{table}[htbp]\centering
\caption{Phase 2 输入输出总览}
\begin{tabular}{@{}llr@{}}
\toprule & \textbf{描述} & \textbf{数量} \\
\midrule
输入 & Phase 1 raw text chunks & 19,464 \\
     & 实际可用（质量筛选后） & 18,597 \\
\midrule
处理 & 采样批次数 & 17 \\
     & 实际处理chunks & 3,206 \\
     & API生成原始QA对 & 4,199 \\
\midrule
输出 & 质量过滤后QA对 & 3,678 \\
     & 训练集（80\%） & 2,941 \\
     & 验证集（10\%） & 366 \\
     & 测试集（10\%） & 371 \\
\bottomrule
\end{tabular}
\end{table}

\section{Prompt 模板设计}

\subsection{五种任务类型}

根据项目书要求，设计了覆盖不同认知层次的五类任务：

\begin{enumerate}[nosep]
    \item \textbf{概念解释}（目标30\%）——训练模型理解和解释领域核心概念
    \item \textbf{章节总结}（目标20\%）——提升长文本理解和结构化总结能力
    \item \textbf{方法对比}（目标20\%）——培养知识迁移和分析比较能力
    \item \textbf{论文摘要理解}（目标15\%）——增强科研文献阅读能力
    \item \textbf{推理分析}（目标15\%）——提高深度分析和安全推理能力
\end{enumerate}

\subsection{系统指令}

所有Prompt共享统一的System Prompt：

\begin{verbatim}
你是一个专业的数据标注助手，负责为计算机科学、密码学和
网络安全领域的教材内容生成高质量的训练样本。
你的回答必须是严格的 JSON 格式，不得包含任何 JSON
之外的文字、解释或 markdown 标记。
每条问答对必须：
1. 问题有明确的知识目标
2. 回答严格依据原文
3. 保留原文的专业术语
4. 如果原文质量不足以生成有效问答对（例如OCR乱码过多），
   直接返回空数组 []
\end{verbatim}

\subsection{模板示例}

\textbf{概念解释模板}——要求生成2-3个概念解释类问答对，问题必须针对文本中的具体概念或术语，回答先给定义再展开说明。

\textbf{方法对比模板}——若文本涉及两种以上方法/算法/方案，设计对比问题，涵盖结构、性能、安全性等维度，可用表格呈现。

\textbf{推理分析模板}——问题具有分析深度，要求解释“为什么”或“如何”，回答包含问题分析、原理推导和解决方案。

完整模板定义位于 \texttt{phase2/templates/prompts.py}。

\section{API批量生成管道}

\subsection{架构}

\begin{verbatim}
Raw Chunk → 质量筛选 → 任务分配 → Prompt构造
    → DeepSeek-Chat API → JSON解析 → 质量验证 → JSONL存储
\end{verbatim}

\subsection{强模型选择}

采用 DeepSeek-Chat 模型（通过 Anthropic SDK 兼容接口调用）作为生成引擎。选择理由：
\begin{itemize}[nosep]
    \item API 兼容 Anthropic Messages 格式，开发成本低
    \item 中英文双语能力均衡，适合处理混合语料
    \item 性价比高，支持大规模批量生成
\end{itemize}

\subsection{任务分配策略}

采用配额驱动的分配算法，确保五种任务类型接近目标比例：

\begin{verbatim}
对每个chunk:
  1. 计算当前每种任务的已完成比例
  2. 找距离目标比例最远的任务类型
  3. 检测chunk内容特征（含"比较"→方法对比加成，
     含"证明/定理"→推理分析加成，
     含"abstract/RFC"→论文摘要加成）
  4. 综合选择最优任务类型
\end{verbatim}

\subsection{批量处理参数}

\begin{table}[htbp]\centering
\caption{批量生成配置}
\begin{tabular}{@{}ll@{}}
\toprule
\textbf{参数} & \textbf{值} \\
\midrule
模型 & deepseek-chat \\
最大输出token & 2,048 \\
温度 & 0.7 \\
输入截断 & 1,500 字符 \\
API调用间隔 & 2 秒 \\
每chunk生成QA对数 & 1--3 对 \\
重试次数 & 3（指数退避） \\
检查点保存间隔 & 每 20 chunks \\
\bottomrule
\end{tabular}
\end{table}

\subsection{错误处理}

实现了三级容错机制：
\begin{enumerate}[nosep]
    \item \textbf{API瞬时故障}：指数退避重试（1s/2s/4s），最多3次
    \item \textbf{JSON解析失败}：正则提取 + 回退到单个对象解析
    \item \textbf{ThinkingBlock处理}：DeepSeek推理模型返回的思考块自动过滤，仅保留TextBlock
\end{enumerate}

整体API成功率为76\%（200 chunks中约152个成功），失败主要原因为chunk内容过短（<150字符）或OCR乱码严重。

\section{数据质量控制}

\subsection{自动过滤规则}

\begin{enumerate}[nosep]
    \item \textbf{长度过滤}：input $<$ 5字符 或 output $<$ 30字符 → 丢弃
    \item \textbf{乱码检测}：有效字符（中英文+数字+标点）占比 $<$ 30\% → 丢弃
    \item \textbf{输出质量}：output $<$ 80字符的浅层回答标记为待审核
\end{enumerate}

\subsection{过滤效果}

\begin{table}[htbp]\centering
\caption{质量过滤统计}
\begin{tabular}{@{}lrl@{}}
\toprule
\textbf{阶段} & \textbf{数量} & \textbf{说明} \\
\midrule
API原始生成 & 4,199 & 17批次全部产出 \\
JSON解析失败 & 38 & API返回格式异常 \\
质量过滤剔除 & 483 & 短回答/乱码/无意义 \\
最终入库 & 3,678 & 通过全部质量检查 \\
\bottomrule
\end{tabular}
\end{table}

\section{数据集划分}

\subsection{分层抽样策略}

按 domain（computer\_science / crypto\_and\_security / comprehensive）进行分层随机抽样，保证每个子集中领域分布一致。划分比例：训练集80\%、验证集10\%、测试集10\%。

\begin{table}[htbp]\centering
\caption{数据集划分结果}
\begin{tabular}{@{}lrrrr@{}}
\toprule
\textbf{领域} & \textbf{总数} & \textbf{训练} & \textbf{验证} & \textbf{测试} \\
\midrule
crypto\_and\_security & 1,843 & 1,474 & 184 & 185 \\
computer\_science      & 1,708 & 1,366 & 170 & 172 \\
comprehensive          & 127   & 101   & 12  & 14 \\
\midrule
\textbf{总计}         & \textbf{3,678} & \textbf{2,941} & \textbf{366} & \textbf{371} \\
\bottomrule
\end{tabular}
\end{table}

\begin{table}[htbp]\centering
\caption{任务类型分布}
\begin{tabular}{@{}lrrrr@{}}
\toprule
\textbf{任务类型} & \textbf{训练集} & \textbf{验证集} & \textbf{测试集} & \textbf{占比} \\
\midrule
概念解释   & 1,455 & 168 & 173 & 49\% \\
推理分析   & 516   & 68  & 62  & 18\% \\
章节总结   & 451   & 65  & 64  & 16\% \\
方法对比   & 336   & 36  & 44  & 11\% \\
论文摘要理解 & 183 & 29  & 28  & 7\% \\
\bottomrule
\end{tabular}
\end{table}

\subsection{输出格式}

每条训练样本为标准 instruction-tuning 格式：

\begin{verbatim}
{
  "instruction": "请解释以下概念\n\n什么是AES加密算法？",
  "output": "AES (Advanced Encryption Standard) 是...",
  "domain": "crypto_and_security",
  "task_type": "概念解释"
}
\end{verbatim}

\subsection{加载方式}

\begin{verbatim}
from datasets import load_dataset
ds = load_dataset("json", data_files={
    "train": "phase2/dataset/train.jsonl",
    "val":   "phase2/dataset/val.jsonl",
    "test":  "phase2/dataset/test.jsonl",
})
\end{verbatim}

\section{长度统计}

\begin{table}[htbp]\centering
\caption{各子集平均长度}
\begin{tabular}{@{}lrr@{}}
\toprule
\textbf{子集} & \textbf{平均Instruction长度} & \textbf{平均Output长度} \\
\midrule
训练集 & 54 字符 & 252 字符 \\
验证集 & 53 字符 & 255 字符 \\
测试集 & 53 字符 & 252 字符 \\
\bottomrule
\end{tabular}
\end{table}

\section{项目脚本清单}

\begin{table}[htbp]\centering
\caption{Phase 2 脚本文件}
\small
\begin{tabular}{@{}ll@{}}
\toprule
\textbf{文件} & \textbf{功能} \\
\midrule
\texttt{templates/prompts.py}        & 5种任务Prompt模板 + System Prompt \\
\texttt{scripts/batch\_generate.py}  & API批量QA生成（任务分配/调用/过滤/存储） \\
\texttt{scripts/split\_dataset.py}   & 分层抽样 + 数据集划分 + HuggingFace格式导出 \\
\texttt{scripts/run\_all\_batches.sh}& 串行批量执行脚本 \\
\texttt{scripts/run\_remaining.sh}   & 补充批次执行脚本 \\
\bottomrule
\end{tabular}
\end{table}

\section{已知局限与改进方向}

\begin{enumerate}[nosep]
    \item \textbf{任务分布偏斜}：概念解释占49\%，超出目标30\%。主要原因是早期批次使用随机权重分配，后改为配额驱动的分配器已修复此问题，但前几批数据已生成。
    \item \textbf{输出长度偏短}：平均output仅252字符。可通过提高Prompt要求（“请给出详细回答，不少于300字”）来改善。
    \item \textbf{论文摘要理解偏少}（7\% vs 目标15\%）：受限于raw chunks中论文/标准类内容比例较低。如需补充，可单独从NIST SP 800系列和RFC文档中生成。
    \item \textbf{comprehensive领域样本不足}（127对）：该领域本身chunk数量有限，跨领域文档仅126个。
\end{enumerate}

\section{第三阶段展望}

第二阶段产出的 \texttt{phase2/dataset/} 目录可直接迁移至GPU设备，进入第三阶段——QLoRA微调：

\begin{enumerate}[nosep]
    \item 环境准备：安装 PyTorch、Transformers、PEFT、BitsAndBytes、Accelerate
    \item 模型加载：Qwen3-4B-Base + 4-bit NF4量化
    \item LoRA配置：rank=32, alpha=64, target\_modules=all-linear
    \item 训练参数：batch\_size=2, grad\_accum=4, lr=2e-4, epochs=3, max\_seq\_len=1024
    \item 预估显存：$\sim$6.1 GB（RTX 4060 8GB可运行）
    \item 评估维度：准确性、完整性、专业性、忠实性
\end{enumerate}

\end{document}
